\begin{center}
{\Large\bfseries Circles}\\[3pt]
{\large\bfseries Test 1}\\[5pt]
Date: \rule{3cm}{0.4pt} \quad Duration: 1 hour 30 minutes \quad Maximum marks: 40
\end{center}
\vspace{-6pt}
\noindent\rule{\linewidth}{0.6pt}\par
\vspace{3pt}
{\centering\small All questions are compulsory. Draw a rough diagram wherever it helps. Show all your steps and give reasons.\par}
\ifteacher
\vspace{4pt}
{\color{thinrule}\noindent\rule{\linewidth}{0.4pt}}\par
\vspace{4pt}
{\small\color{muted}\linespread{1.15}\selectfont
\textbf{Paper design.} The paper follows the CBSE pattern. Section A: $6\times1$ mark (4 multiple-choice questions, 1 fill-in-the-blank question and 1 assertion--reason question). Section B: $4\times2$ marks. Section C: $4\times3$ marks. Section D: $2\times5$ marks. Section E: 1 case-based question of 4 marks. Total: 40 marks.\nl
\textbf{Theorem numbers.} The paper uses the theorem numbers of the class notes (Circles classes 1--4). Theorems 1--4 carry about 27 marks. Theorems 5--8 carry about 13 marks. Theorems 9--12 are not assessed.\nl
\textbf{Competency focus.} Question 2 tests direct use of a result. The other questions are competency-based. Students find an error (Q7), work backwards (Q10, Q16), write proofs (Q12, Q15), construct a figure (Q13), compare chords (Q5, Q14) and apply results to real situations (Q8, Q13, Q16, Q17).\nl
\textbf{Concepts assessed.} Equal chords and angles at the centre, Theorems 1 and 2 (Q1, 8, 12, 17); the line from the centre to a chord, Theorems 3 and 4 (Q2, 3, 6, 7, 9, 10, 11, 14, 15, 16, 17); the circle through three points, Theorem 5 (Q4, 13); chords and their distance from the centre, Theorems 6--8 (Q5, 9, 14, 16, 17).\nl
\textbf{Common misconceptions.} Joining the centre to any point of a chord and taking the line as perpendicular (Q7); using the full chord in place of half the chord in Pythagoras theorem (Q2, 10, 11); adding the distances in every case of parallel chords (Q11).\nl
\textbf{Marking.} Award method marks for a correct set-up, even where the final computation contains an arithmetic error. In proofs, award marks for the correct congruence rule with all three conditions and their reasons.\par}
\vspace{4pt}
{\color{thinrule}\noindent\rule{\linewidth}{0.4pt}}\par
\vspace{3pt}
\fi
